\section{A Structured State for Budget Decisions}

\subsection{One Proof Agent, Two Execution Regimes}

The system exposes two explicit execution modes. The \emph{minimal} mode is a
linear refinement loop:
\[
\textsc{Propose}\rightarrow\textsc{Check}\rightarrow
\textsc{Remember}\rightarrow\textsc{Refine}.
\]
It carries compact memory of established facts, failed approaches, Lean API
lessons, and open goals. This mode is both a practical solver for local
theorems and the paper's main agentic baseline.

The \emph{structured} mode is the substrate for budget-aware control. It uses
the same proof model, checker, retrieval interface, and safety review, but
represents several live obligations and alternatives explicitly. A run chooses
one mode at startup; it never silently falls back from structured search to
the minimal loop. This keeps measurements interpretable.

We deliberately use one model role, the \emph{Proof Agent}. Mathematical
planning, Lean implementation, and error interpretation remain in the same
reasoning context because a Lean error can expose a missing mathematical step,
a library mismatch, a poor representation, or a merely local tactic defect.
The controller schedules computations but does not assign failures to a chain
of role-specialized agents.

\subsection{Proof Workspace}

A workspace \(W\) is an immutable, versioned state with five principal layers.

\paragraph{Formal specification.}
The original problem, fixed Lean statement, imports, and proof hole define the
statement that all accepted candidates must preserve.

\paragraph{Obligation graph.}
Proof obligations are nodes in a directed acyclic graph. An edge
\(o\rightarrow d\) states that obligation \(o\) requires dependency \(d\).
All dependencies of the root must be accepted, so the graph expresses the AND
structure of a complex proof. A statement, assumption, or dependency change
creates a new obligation version; old versions remain as provenance but cannot
be used as current facts.

\paragraph{Proof branches.}
Each obligation may have several branches representing alternative arguments
or Lean realizations, giving OR structure. A branch stores its parent, argument
steps, Lean artifact, alignment links, observations, failure hypotheses, last
action, progress signal, and status. A failed Lean implementation does not
automatically refute the mathematical strategy that produced it.

\paragraph{Evidence and verified facts.}
Checker diagnostics become append-only observations with source spans,
declaration identifiers, goal fingerprints, and raw trace references.
Accepted facts are registered only after both Lean checking and safety review,
and they retain the obligation version and artifact provenance. Capability
audits record narrower mechanical facts such as whether an identifier exists
or a minimal signature elaborates; they do not prove the target theorem.

\paragraph{Argument--Lean alignment.}
Natural-language argument steps link to Lean declarations, source spans, or
goals with one of three relations: implemented, partial, or unaligned. This
alignment lets the controller distinguish a local implementation repair from
an action that changes mathematical content. When evidence is insufficient,
the state records uncertainty instead of inventing a unique cause.

\subsection{Typed Actions and Deterministic Reduction}

Every proposed action contains a kind, target branch, affected argument steps,
rationale, and an upper bound on allowed mutations. For example, an
implementation repair may change a Lean artifact and its alignment but may not
silently revise the theorem assumptions. Structural actions such as
decomposition and representation change have separate payloads and transition
rules.

Tool outcomes are folded into the workspace by deterministic reducers. An
accepted implementation stores an artifact and verified fact; a rejected
candidate appends checker observations; decomposition supersedes the old parent
version and creates helper obligations; a missing mechanical capability blocks
the affected dependency closure; and new verified evidence may reactivate a
dormant branch. Because reducers create successors rather than mutating old
states, every policy decision can be replayed from the trace.

\subsection{Dependency-Aware Frontier and Assembly}

A branch is schedulable only if it is active, its obligation is open, and all
of that obligation's dependencies are accepted. Readiness is a gate rather
than a small priority penalty: repeatedly selecting an unready parent wastes
budget without changing the proof state. Within the ready set, the frontier
uses a stable baseline order before the budget-aware score is applied.

When every active obligation has a compatible accepted branch, a solution
tracker selects one artifact per obligation. Helper declarations are injected
before the root theorem and the entire source is checked again. The controller
reserves one verifier call for this final assembly. A local collection of
accepted artifacts is not reported as a solved theorem unless the integrated
file passes both the checker and statement-preservation/anti-cheating review.

\subsection{Context Projection}

The full workspace may be larger than a useful prompt. For a chosen branch we
project the root goal, current obligation and version, dependency closure,
accepted facts, aligned argument steps, deduplicated recent observations,
failure hypotheses, and compact sibling summaries. This projection is the
single source of prompt context; a separate summarizer may compress it but may
not create a parallel set of ``facts.'' The projection cost itself is included
in input-token accounting.
